\documentclass[12pt,a4paper]{article}
\usepackage[utf8]{vietnam}
\usepackage{amsmath,amssymb}
\usepackage{geometry}
\geometry{
  top=12mm,
  bottom=12mm,
  left=14mm,
  right=14mm,
  footskip=6mm,
  headheight=14pt
}

\usepackage{tikz}
\usetikzlibrary{arrows.meta,calc,shapes.geometric}
\usepackage[most]{tcolorbox}
\usepackage{fontawesome}
\usepackage{fancyhdr}
\usepackage{enumitem}

\pagestyle{fancy}
\fancyhf{}
\fancyhead[L]{\footnotesize\color{teal!80!black}\faBook\ \textbf{SỔ TAY KIẾN THỨC TOÁN 3}}
\fancyhead[R]{\footnotesize\color{teal!80!black}\textbf{Lớp Toán Thầy Lĩnh -- loptoanthaylinh.vn}}
\fancyfoot[C]{\footnotesize\color{gray}-- Sổ tay kiến thức: Bài 10. Bảng nhân 7, bảng chia 7 --}
\renewcommand{\headrulewidth}{0.4pt}
\renewcommand{\footrulewidth}{0pt}

\setlength{\parindent}{0pt}
\setlength{\parskip}{0pt}

\begin{document}

% --- BANNER BÀI HỌC ---
\begin{tcolorbox}[
  enhanced,
  colback=teal!8!white,
  colframe=teal!70!black,
  arc=3mm, boxrule=1.2pt,
  left=3mm, right=3mm, top=2mm, bottom=2mm,
  halign=center
]
  {\LARGE\bfseries\color{teal!90!black}\faBookmark\ BÀI 10: BẢNG NHÂN 7, BẢNG CHIA 7}\\[1mm]
  {\itshape\footnotesize\color{black!75} Quy tắc đếm thêm 7, các mốc liên hệ thực tế 1 tuần có 7 ngày trên 1 trang A4}
\end{tcolorbox}

\vspace{-2mm}

\noindent
\begin{minipage}[t]{0.49\textwidth}\vspace{0pt}
\begin{tcolorbox}[
  enhanced,
  colback=blue!4!white,
  colframe=blue!75!black,
  arc=2mm, boxrule=0.8pt,
  left=2mm, right=2mm, top=1.8mm, bottom=1.8mm,
  title={\footnotesize\bfseries\color{white}\faTimes\ 1. Bảng nhân 7 (Đếm thêm 7)}
]
{\footnotesize
\begin{itemize}\setlength{\itemsep}{0.6pt}\setlength{\leftskip}{-2.5mm}
  \item Dãy tích của bảng 7 tăng đều 7 đơn vị:
  $$7, 14, 21, 28, 35, 42, 49, 56, 63, 70$$
  \item \textbf{Liên hệ 1 tuần lễ có 7 ngày:}
  \begin{itemize}\setlength{\itemsep}{0.5pt}\setlength{\leftskip}{-2.5mm}
    \item 2 tuần: $7 \times 2 = 14$ ngày; 3 tuần: $7 \times 3 = 21$ ngày.
    \item 4 tuần: $7 \times 4 = 28$ ngày; 5 tuần: $7 \times 5 = 35$ ngày.
  \end{itemize}
  \item Mốc đặc biệt: $7 \times 7 = \mathbf{49}$; $7 \times 8 = \mathbf{56}$; $7 \times 9 = \mathbf{63}$.
\end{itemize}
}
\end{tcolorbox}
\end{minipage}
\hfill
\begin{minipage}[t]{0.49\textwidth}\vspace{0pt}
\begin{tcolorbox}[
  enhanced,
  colback=teal!4!white,
  colframe=teal!75!black,
  arc=2mm, boxrule=0.8pt,
  left=2mm, right=2mm, top=1.8mm, bottom=1.8mm,
  title={\footnotesize\bfseries\color{white}\faShareAlt\ 2. Bảng chia 7 (Từ nhân suy ra chia)}
]
{\footnotesize
\begin{itemize}\setlength{\itemsep}{0.6pt}\setlength{\leftskip}{-2.5mm}
  \item Lấy tích chia cho 7 thì được thừa số kia:
  \begin{tabular}{ll}
    $7 : 7 = 1$ & $42 : 7 = 6$\\
    $14 : 7 = 2$ & $49 : 7 = 7$\\
    $21 : 7 = 3$ & $56 : 7 = 8$\\
    $28 : 7 = 4$ & $63 : 7 = 9$\\
    $35 : 7 = 5$ & $70 : 7 = 10$
  \end{tabular}
  \item Muốn tính $56 : 7$, ta nhẩm: "Mấy nhân $7$ bằng $56$?". Vì $7 \times 8 = 56 \implies 56 : 7 = 8$.
\end{itemize}
}
\end{tcolorbox}
\end{minipage}

\vspace{-2.5mm}
% --- KHỐI 3: SƠ ĐỒ TƯ DUY TIKZ ---
\noindent
\begin{tcolorbox}[
  enhanced,
  colback=orange!4!white,
  colframe=orange!80!black,
  arc=2mm, boxrule=0.8pt,
  left=2mm, right=2mm, top=1.5mm, bottom=1.5mm,
  title={\footnotesize\bfseries\color{white}\faSitemap\ SƠ ĐỒ TƯ DUY: CẦU NỐI TUẦN LỄ \& BẢNG NHÂN 7}
]
\centering
\begin{tikzpicture}[
  box/.style={draw, rounded corners=2pt, align=center, font=\scriptsize\bfseries, inner sep=3pt, minimum height=0.75cm},
  arr/.style={-{Stealth[length=2mm]}, thick}
]
  \node[box, fill=teal!15, draw=teal!80!black] (w1) at (0, 0) {1 TUẦN\\$7\text{ ngày}$};
  \node[box, fill=blue!12, draw=blue!80!black] (w2) at (2.4, 0) {3 TUẦN\\$7 \times 3 = \mathbf{21}$};
  \node[box, fill=purple!15, draw=purple!80!black] (w3) at (5.0, 0) {5 TUẦN\\$7 \times 5 = \mathbf{35}$};
  \node[box, fill=orange!20, draw=orange!80!black] (w4) at (7.6, 0) {7 TUẦN\\$7 \times 7 = \mathbf{49}$};
  \node[box, fill=red!15, draw=red!80!black] (w5) at (10.0, 0) {9 TUẦN\\$7 \times 9 = \mathbf{63}$};

  \draw[arr, teal!80!black] (w1) -- (w2);
  \draw[arr, blue!80!black] (w2) -- (w3);
  \draw[arr, purple!80!black] (w3) -- (w4);
  \draw[arr, orange!80!black] (w4) -- (w5);
\end{tikzpicture}
\end{tcolorbox}

\vspace{-2.5mm}
% --- KHỐI 4: VÍ DỤ & MẸO TRÁNH BẪY ---
\noindent
\begin{minipage}[t]{0.49\textwidth}\vspace{0pt}
\begin{tcolorbox}[
  enhanced,
  colback=green!4!white,
  colframe=green!60!black,
  arc=2mm, boxrule=0.8pt,
  left=2mm, right=2mm, top=1.8mm, bottom=1.8mm,
  title={\footnotesize\bfseries\color{white}\faCheckCircle\ Ví dụ thực hành chuẩn mực}
]
{\footnotesize
\textbf{1. Tính nhẩm:}\\
$7 \times 4 = \mathbf{28}$;\quad $7 \times 8 = \mathbf{56}$.\\
$49 : 7 = \mathbf{7}$;\quad $63 : 7 = \mathbf{9}$.\\
\textbf{2. Đổi tuần sang ngày:}\\
$6$ tuần lễ có số ngày là: $7 \times 6 = \mathbf{42}$ (ngày).\\
\textbf{3. Bài toán thực tế:}\\
Có $35$ bông hoa cắm đều vào $7$ lọ. Mỗi lọ có: $35 : 7 = \mathbf{5}$ (bông hoa).
}
\end{tcolorbox}
\end{minipage}
\hfill
\begin{minipage}[t]{0.49\textwidth}\vspace{0pt}
\begin{tcolorbox}[
  enhanced,
  colback=red!4!white,
  colframe=red!70!black,
  arc=2mm, boxrule=0.8pt,
  left=2mm, right=2mm, top=1.8mm, bottom=1.8mm,
  title={\footnotesize\bfseries\color{white}\faExclamationTriangle\ Mẹo tránh bẫy trắc nghiệm}
]
{\footnotesize
\begin{itemize}\setlength{\itemsep}{1.5pt}\setlength{\leftskip}{-2.5mm}
  \item \textbf{Cặp số 56 7 8 thần thánh:}\\
  $5, 6, 7, 8$ tạo thành chuỗi liên tiếp: $\mathbf{56 = 7 \times 8}$. Câu thần chú giúp không bao giờ nhầm $7 \times 8$!
  \item \textbf{Bẫy hay nhầm $7 \times 6$ và $7 \times 7$:}\\
  $7 \times 6 = 42$, còn $7 \times 7 = 49$. Tránh nhầm $7 \times 6 = 48$ của bảng 6!
\end{itemize}
}
\end{tcolorbox}
\end{minipage}


\end{document}
